\subsection{Parameter Integrals and Differentiation}

\begin{corollary} \label{cor:dominated-parameter-continuity} \hlt{Continuity Under a Lebesgue Integral}\\
Let $f_t(x)\defeq f(x,t)$ be measurable $\forall t\in[a,b]$, and suppose $\forall t\in[a,b],\ \abs{f_t}\leq g\in\mathscr{L}^1(\mu)$. Then:
\begin{enumerate}[label=(\roman*)]
\setlength{\itemsep}{0pt}
\item If $t_0\in[a,b]$ and $f_t\to f_{t_0}$ a.e. as $t\to t_0$, then $\lim\limits_{t\to t_0}\int_X f_t\,d\mu=\int_X f_{t_0}\,d\mu$.
\item If $t\mapsto f(x,t)$ is continuous on $[a,b]$ for a.e. $x$, then $F(t)\defeq\int_X f_t\,d\mu$ is continuous on $[a,b]$.
\end{enumerate}
\begin{proof}
\begin{enumerate}[label=(\roman*)]
\setlength{\itemsep}{0pt}
\item For any $t_n\to t_0$, set $u_n\defeq f_{t_n}$. Then $u_n\to f_{t_0}$ a.e. and $\abs{u_n}\leq g$, so dominated convergence gives $\int_Xu_n\,d\mu\to\int_Xf_{t_0}\,d\mu$. The sequential criterion gives the claim.
\item If $t\mapsto f(x,t)$ is continuous a.e., then $\forall t_0\in[a,b]$, $f_t\to f_{t_0}$ a.e. as $t\to t_0$.\\
By item (i), $\forall t_0\in[a,b]$, $F(t)\to F(t_0)$, so $F$ is continuous on $[a,b]$.
\end{enumerate}
\hlt{Structural Commentary}: Parameter continuity reduces to dominated convergence along sequences.\\
\hlt{Key Insight}: The parameter disappears once $t_n$ is frozen into a sequence of functions.
\end{proof}
\end{corollary}

\begin{theorem} \label{thm:differentiation-under-lebesgue-integral} \hlt{Differentiation Under a Lebesgue Integral}\\
Let $f_t(x)\defeq f(x,t)$ be measurable $\forall t\in[a,b]$, and suppose $f_{t_0}\in\mathscr{L}^1(\mu)$ for some $t_0\in[a,b]$.\\
Assume $t\mapsto f(x,t)$ is continuous on $[a,b]$ and differentiable on $(a,b)$ $\forall x\in X$, and $\exists g\in\mathscr{L}^1(\mu)$ such that $\forall(x,t)\in X\times(a,b)$, $\abs{\partial_tf(x,t)}\leq g(x)$. For $t\in(a,b)$, set $D_t(x)\defeq\partial_tf(x,t)$; set $F(t)\defeq\int_Xf_t\,d\mu$. Then:
\begin{enumerate}[label=(\roman*)]
\setlength{\itemsep}{0pt}
\item $\forall t\in(a,b)$, $D_t$ is measurable.
\item $\forall t\in[a,b]$, $f_t\in\mathscr{L}^1(\mu)$.
\item $F$ is differentiable on $(a,b)$, and $\forall t\in(a,b)$, $F'(t)=\int_XD_t\,d\mu$.
\end{enumerate}
\begin{proof}
\begin{enumerate}[label=(\roman*)]
\setlength{\itemsep}{0pt}
\item Fix $t\in(a,b)$ and choose $h_k\to0$ with $t+h_k\in[a,b]$, $h_k\neq0$.\\
The quotients $\frac{f_{t+h_k}-f_t}{h_k}$ are measurable and converge pointwise to $D_t$, so $D_t$ is measurable.
\item By the mean value theorem, $\abs{f_t-f_{t_0}}\leq \abs{t-t_0}g$, so $\abs{f_t}\leq\abs{f_{t_0}}+\abs{t-t_0}g\in\mathscr{L}^1(\mu)$.
\item If $t_n\to t\in(a,b)$ and $t_n\neq t$, then $\frac{F(t_n)-F(t)}{t_n-t}=\int_X\frac{f_{t_n}-f_t}{t_n-t}\,d\mu$. The integrand converges pointwise to $D_t$ and has modulus at most $g$, so dominated convergence gives $\frac{F(t_n)-F(t)}{t_n-t}\to\int_XD_t\,d\mu$.
\end{enumerate}
\hlt{Structural Commentary}: Bounded parameter derivatives dominate the difference quotients.\\
\hlt{Key Insight}: Differentiating the integral is just applying dominated convergence to secant slopes.
\end{proof}
\end{theorem}

\begin{corollary} \label{cor:single-point-differentiation-under-lebesgue-integral} \hlt{A Pointwise Difference-Quotient Test}\\
Let $f_t(x)\defeq f(x,t)$ be measurable $\forall t\in[a,b]$, let $f_{t_0}\in\mathscr{L}^1(\mu)$ for some $t_0\in[a,b]$, and fix $t_1\in[a,b]$.\\
Assume $\partial_tf(x,t_1)$ exists $\forall x$. For $t\neq t_1$, set $Q_t(x)\defeq\frac{f_t(x)-f_{t_1}(x)}{t-t_1}$, $\abs{Q_t}\leq g$ for some $g\in\mathscr{L}^1(\mu)$.\\
Set $F(t)\defeq\int_Xf_t\,d\mu$. Then:
\begin{enumerate}[label=(\roman*)]
\setlength{\itemsep}{0pt}
\item $\forall t\in[a,b]$, $f_t\in\mathscr{L}^1(\mu)$.
\item $F$ is differentiable at $t_1$, and $F'(t_1)=\int_X\partial_tf(x,t_1)\,d\mu$.
\end{enumerate}
\begin{proof}
\begin{enumerate}[label=(\roman*)]
\setlength{\itemsep}{0pt}
\item If $t_0\neq t_1$, then $\abs{f_{t_1}}\leq\abs{f_{t_0}}+\abs{t_0-t_1}g$; if $t_0=t_1$, this is given.\\
Hence $f_{t_1}\in\mathscr{L}^1(\mu)$. For $t\neq t_1$, $\abs{f_t}\leq\abs{f_{t_1}}+\abs{t-t_1}g$, so $f_t\in\mathscr{L}^1(\mu)$.
\item For any $t_n\to t_1$ with $t_n\neq t_1$, $Q_{t_n}\to\partial_tf(\cdot,t_1)$ pointwise and $\abs{Q_{t_n}}\leq g$.\\
Dominated convergence gives $\frac{F(t_n)-F(t_1)}{t_n-t_1}\to\int_X\partial_tf(x,t_1)\,d\mu$.
\end{enumerate}
\hlt{Structural Commentary}: One parameter point needs only domination of the secants into that point.\\
\hlt{Key Insight}: A global derivative bound is sufficient, but not necessary, for a one-point formula.
\end{proof}
\end{corollary}

\begin{example} \label{ex:differentiating-under-integral-dominated-convergence} \hlt{Differentiating Under the Integral Sign}\\
For $t>0$, set $F(t)\defeq\int_0^\infty x^2e^{-tx}\,dx$. For $t_0>0$ and $0<\abs{h}<\frac{t_0}{2}$, put $Q_h(x)\defeq\frac{x^2e^{-(t_0+h)x}-x^2e^{-t_0x}}{h}$.\\
Then $Q_h\to -x^3e^{-t_0x}$ pointwise and $\abs{Q_h}\leq x^3e^{-\frac{t_0x}{2}}\in\mathscr{L}^1(0,\infty)$. Hence $F'(t_0)=-\int_0^\infty x^3e^{-t_0x}\,dx$.\\
\hlt{Structural Commentary}: Parameter differentiation is a limit-interchange problem.\\
\hlt{Key Insight}: Prove a local-in-parameter integrable majorant for the difference quotients.
\end{example}

\begin{example} \label{ex:laplace-transform-powers-factorials} \hlt{Laplace Transform Powers and Factorials}\\
For $t>0$, set $I(t)\defeq\int_0^\infty e^{-tx}\,dx$. Then $I(t)=\lim\limits_{b\to+\infty}\frac{1-e^{-tb}}{t}=t^{-1}$.\\
For $[a,b]\subset(0,+\infty)$ and $k\geq0$, $0\leq x^ke^{-tx}\leq x^ke^{-ax}\in\mathscr{L}^1(0,\infty)$, so repeated differentiation under the integral gives $I^{(n)}(t)=(-1)^n\int_0^\infty x^ne^{-tx}\,dx$. Since $I^{(n)}(t)=(-1)^nn!t^{-n-1}$, $\forall n\in\N\cup\{0\}$, $\int_0^\infty x^ne^{-x}\,dx=n!$.\\
\hlt{Structural Commentary}: A parameter integral turns powers of $x$ into derivatives in $t$.\\
\hlt{Key Insight}: The exponential envelope $e^{-ax}$ controls all nearby parameter values $t\geq a$.
\end{example}

\begin{corollary} \label{cor:riemann-parameter-integral-interchange} \hlt{One Lebesgue and One Riemann Integral}\\
Let $f_t(x)\defeq f(x,t)$, $t\in[a,b]$, with $f_t$ measurable, $t\mapsto f(x,t)$ continuous a.e., and $\forall t\in[a,b],\ \abs{f_t}\leq g\in\mathscr{L}^1(\mu)$.\\
Then, with Riemann integrals in $t$, $\int_a^b\left[\int_X f_t\,d\mu\right]dt =\int_X\left[\int_a^b f(x,t)\,dt\right]d\mu(x)$.
\begin{proof}
Set $h_t(x)\defeq\int_a^t f(x,s)\,ds$. Then $h_t$ is measurable, $\abs{h_t}\leq(b-a)g$, and $\partial_th_t=f_t$.\\
By differentiation under the integral, $H(t)\defeq\int_Xh_t\,d\mu$ satisfies $H'(t)=\int_Xf_t\,d\mu$.\\
Hence $\int_a^b\left[\int_Xf_t\,d\mu\right]dt=H(b)-H(a)=\int_X\left[\int_a^b f(x,t)\,dt\right]d\mu(x)$.\\
\hlt{Structural Commentary}: Finite parameter integrals move through by differentiating primitives.\\
\hlt{Key Insight}: Integrate first in the parameter up to $t$, then differentiate the resulting Lebesgue integral.
\end{proof}
\end{corollary}

\begin{corollary} \label{cor:improper-riemann-parameter-integral-interchange} \hlt{Improper Parameter Interchange}\\
Let $f_t(x)\defeq f(x,t)$ be measurable $\forall t\in\R$, and let $t\mapsto f(x,t)$ be continuous for a.e. $x$.\\
Suppose $\exists g,h\in\mathscr{L}^1(\mu)$ such that $\abs{f_t}\leq g$ $\forall t$ and $\int_\R\abs{f(x,t)}\,dt\leq h(x)$ for a.e. $x$.\\
Put $G(t)\defeq\int_Xf_t\,d\mu$, $H(x)\defeq\int_\R f(x,t)\,dt$, and $H_R(x)\defeq\int_{-R}^{R}f(x,t)\,dt$. Then:
\begin{enumerate}[label=(\roman*)]
\setlength{\itemsep}{0pt}
\item $H\in\mathscr{L}^1(\mu)$.
\item $\lim\limits_{R\to+\infty}\int_{-R}^{R}G(t)\,dt=\int_XH\,d\mu$.
\end{enumerate}
Thus $\int_\R G\defeq\lim\limits_{R\to+\infty}\int_{-R}^{R}G(t)\,dt$ exists and equals $\int_XH\,d\mu$.
\begin{proof}
\begin{enumerate}[label=(\roman*)]
\setlength{\itemsep}{0pt}
\item Since $\abs H\leq\int_\R\abs{f(x,t)}\,dt\leq h$ a.e. and $h\in\mathscr{L}^1(\mu)$, $H\in\mathscr{L}^1(\mu)$ by Theorem \ref{thm:absolute-integrability-and-domination}(ii).
\item For $R>0$, finite-interval interchange gives $\int_{-R}^{R}G(t)\,dt=\int_XH_R\,d\mu$.\\
Since $H_R\to H$ a.e. and $\abs{H_R}\leq h$, dominated convergence gives $\int_XH_R\,d\mu\to\int_XH\,d\mu$.
\end{enumerate}
\hlt{Structural Commentary}: Finite interchanges pass to improper limits under an $L^1$ bound.\\
\hlt{Key Insight}: Truncate in $t$, interchange on finite intervals, then pass $R\to+\infty$ by domination in $x$.
\end{proof}
\end{corollary}
